%%%%%%%%%%%%%%%%%%%%%%%%%%%%%%%%%%%%%%%%%%%%%%%%%%%%%%%%%%%%%%%%%
% Visualization & Description Table
%%%%%%%%%%%%%%%%%%%%%%%%%%%%%%%%%%%%%%%%%%%%%%%%%%%%%%%%%%%%%%%%%
\bgroup
\begin{table*}
\caption{Example visualizations and descriptions from our corpus. Paragraph breaks in rows \texttt{A} and \texttt{B} indicate a description authored by a unique participant from our corpus gathering survey~\sectionlabel{section:study1design}, while row \texttt{C} shows an curated exemplar description from our evaluation~\sectionlabel{section:study2design}.
% We include the entire corpus, and all of the curated exemplars used in the evaluation, in the supplementary material.
}
\centering
\def\arraystretch{1.5}

\begin{tabular}{m{0.2cm} m{7.5cm} m{9cm}}
\hline
\multicolumn{2}{l}{\textsc{visualization}} & \textsc{description}
\\
\hline

\cellcolor{GrayLight}\textbf{\texttt{A}} &
\includegraphics[scale=0.33]{figures/1134_bar_easy_journalism.png}
\textbf{\texttt{[bar, easy, journalism]}}
&
% \cellcolor{GrayLight}
{\sffamily\smaller

    [1]~\hlc[Level1]{This is a vertical bar chart entitled ``COVID-19 mortality rate by age'' that plots Mortality rate by Age. Mortality rate is plotted on the vertical y-axis from 0 to 15\%. Age is plotted on the horizontal x-axis in bins: 10-19, 20-29, 30-39, 40-49, 50-59, 60-69, 70-79, 80+.}
    [2]~\hlc[Level2]{The highest COVID-19 mortality rate is in the 80+ age range, while the lowest mortality rate is in 10-19, 20-29, 30-39, sharing the same rate.}
    [3]~\hlc[Level3]{COVID-19 mortality rate does not linearly correspond to the demographic age.}

    [4]~\hlc[Level3]{The mortality rate increases with age, especially around 40-49 years and upwards.}
    [5]~\textcolor{white}{\hlc[Level4]{This relates to people's decrease in their immunity and the increase of co-morbidity with age.}}

    [6]~\hlc[Level3]{The mortality rate increases exponentially with older people.}
    [7]~\hlc[Level2]{There is no difference in the mortality rate in the range between the age of 10 and 39.}
    [8]~\hlc[Level3]{The range of ages between 60 and 80+ are more affected by COVID-19.}

    [9]~\textcolor{white}{\hlc[Level4]{We can observe that the mortality rate is higher starting at 50 years old due to many complications prior.}}
    [10]~\hlc[Level3]{As we decrease the age, we also decrease the values in mortality by a lot, almost to none.}
}
\\
\hline

\cellcolor{GrayLight}\textbf{\texttt{B}} &
\includegraphics[scale=0.53]{figures/2220_line_medium_business.png}
\textbf{\texttt{[line, medium, business]}}
& 
% \cellcolor{GrayLight}
{\sffamily\smaller

    [1]~\hlc[Level1]{This is a multi-line chart entitled ``Big Tech Stock Prices'' that plots price by date. The corporations include AAPL (Apple), AMZN (Amazon), GOOG (Google), IBM (IBM), and MSFT (Microsoft). The years are plotted on the horizontal x-axis from 2000 to 2010 with an increment of 2 years. The prices are plotted on the vertical y-axis from 0 to 800 with an increment of 200.}
    [2]~\hlc[Level2]{GOOG has the greatest price over time. MSFT has the lowest price over time.}
    [3]~\hlc[Level3]{Prices of particular Big Tech corporations seem to fluctuate but nevertheless increase over time. Years 2008-2009 are exceptions as we can see an extreme drop in prices of all given corporations.}
    [4]~\textcolor{white}{\hlc[Level4]{The big drop in prices was caused by financial crisis of 2007-2008. The crisis culminated with the bankruptcy of Lehman Brothers on September 15, 2008 and an international banking crisis.}}
    
    [5]~\textcolor{white}{\hlc[Level4]{At the beginning of 2008, every of this stock price went down, likely due to the financial crisis.}}
    [6]~\hlc[Level3]{Then they have risen again and dropped again, more so than previously.}
    
    [7]~\hlc[Level2]{GOOG has the highest price over the years. MSFT has the lowest price over the years.}
    [8]~\hlc[Level3]{GOOG quickly became the richest one of the Big Tech corporations.}
    [9]~\textcolor{white}{\hlc[Level4]{GOOG had experienced some kind of a crisis in 2009, because their prices drop rapidly, but then rebounded.}}
}
\\
\hline

\cellcolor{GrayLight}\textbf{\texttt{C}} &
\includegraphics[scale=0.82]{figures/3310_scatter_hard_academic.png}
\textbf{\texttt{[scatter, hard, academic]}}
&
% \cellcolor{GrayLight}
{\sffamily\smaller

    [1]~\hlc[Level1]{This is a scatter plot entitled ``Born in 2016: Life Expectancy Gap by Gender and Income'' that plots Women Life Expectancy at Birth (Years) by Men Life Expectancy at Birth (Years). The Women Life Expectancy at Birth is plotted on the vertical y-axis from 40 to 90 years. The Men Life Expectancy at Birth is plotted on the horizontal x-axis from 40 to 90 years. High Income Countries are plotted in dark green. Low Income Countries are plotted in light green. A 45 degree line from the origin represents Equal Life Expectancy.}
    [2]~\hlc[Level2]{For low income countries, the average life expectancy is 60 years for men and 65 years for women. For high income countries, the average life expectancy is 77 years for men and 82 years for women.}
    [3]~\hlc[Level3]{Overall, women have a slightly higher life expectancy than men. Women live around 5 to 10 years longer than men. The low income countries are more scattered than the high income countries. There is a visible gap between high and low income countries, indicated by the Income-Age Divide line.}
    [4]~\textcolor{white}{\hlc[Level4]{People living in low-income countries tend to have a lower life expectancy than the people living in high-income countries, likely due to many societal factors, including access to healthcare, food, other resources, and overall quality of life. People who live in lower income countries are more likely to experience deprivation and poverty, which can cause related health problems.}}
}
\\
\hline
\end{tabular}
\label{table:visualizations}
\end{table*}
%%%%%%%%%%%%%%%%%%%%%%%%%%%%%%%%%%%%%%%%%%%%%%%%%%%%%%%%%%%%%%%%%
